% ==============================================================================
% SECCIÓN 04.01: FUNCIÓN DE DENSIDAD DE PROBABILIDAD (PDF) Y SOPORTE CONTINUO
% ==============================================================================

\subsection*{Enunciados de los problemas --- Sección 04.01}

\subsubsection*{Nivel Fundamental}

\begin{problema}
	\label{prob:4.1.1}
	Considere la función $f(x) = c \cdot e^{-x}$ para $x \ge 0$ y $f(x) = 0$ para $x < 0$, donde $c$ es una constante positiva.
	\begin{enumerate}
		\item Determine el valor de $c$ tal que $f(x)$ sea una función de densidad de probabilidad (PDF) válida.
		\item Calcule la probabilidad $P(1 \le X \le 3)$ integrando la PDF en el intervalo.
		\item Calcule la función de distribución acumulada (CDF) $F(x)$ para todo $x \in \mathbb{R}$.
	\end{enumerate}
\end{problema}

\begin{problema}
	\label{prob:4.1.2}
	Una variable aleatoria continua $X$ tiene PDF dada por:
	$$f(x) = \begin{cases} c \cdot (4x - 2x^{2}) & \text{si } 0 \le x \le 2 \\ 0 & \text{en otro caso} \end{cases}$$
	\begin{enumerate}
		\item Encuentre la constante de normalización $c$ tal que $\int_{-\infty}^{\infty} f(x)\,dx = 1$.
		\item Calcule $P(X \ge 1)$ y $P(0.5 \le X \le 1.5)$.
	\end{enumerate}
\end{problema}

\begin{problema}
	\label{prob:4.1.3}
	Sea $X$ una variable aleatoria con PDF triangular:
	$$f(x) = \begin{cases} x & \text{si } 0 \le x \le 1 \\ 2 - x & \text{si } 1 < x \le 2 \\ 0 & \text{en otro caso} \end{cases}$$
	\begin{enumerate}
		\item Verifique que $\int_{-\infty}^{\infty} f(x)\,dx = 1$.
		\item Calcule la moda (el valor donde $f$ es máxima) y la mediana (el valor $m$ tal que $P(X \le m) = 0.5$).
	\end{enumerate}
\end{problema}

\subsubsection*{Nivel Operativo}

\begin{problema}
	\label{prob:4.1.4}
	Sea $X$ una variable aleatoria continua con CDF $F(x) = 1 - e^{-x^{2}/2}$ para $x \ge 0$ y $F(x) = 0$ para $x < 0$ (distribución de Rayleigh).
	\begin{enumerate}
		\item Derive la PDF $f(x)$ aplicando $f(x) = F'(x)$.
		\item Verifique que $\int_{0}^{\infty} f(x)\,dx = 1$ usando la sustitución $u = x^{2}/2$.
		\item Calcule la mediana $m$ tal que $F(m) = 0.5$.
	\end{enumerate}
\end{problema}

\begin{problema}
	\label{prob:4.1.5}
	Considere la PDF definida por:
	$$f(x) = \begin{cases} \dfrac{3}{8}(2 - x^{2}) & \text{si } -1 \le x \le 1 \\ 0 & \text{en otro caso} \end{cases}$$
	\begin{enumerate}
		\item Verifique que $f$ integra a $1$.
		\item Calcule $P(-0.5 \le X \le 0.5)$ y $P(X \ge 0)$.
		\item Halle la CDF $F(x)$ para $x \in [-1, 1]$ y grafique cualitativamente su forma.
	\end{enumerate}
\end{problema}

\begin{problema}
	\label{prob:4.1.6}
	En un sistema de comunicaciones, el tiempo de espera $X$ (en milisegundos) hasta la recepción de un paquete tiene PDF exponencial con parámetro $\lambda = 0.5$ ms$^{-1}$:
	$$f(x) = 0.5 \cdot e^{-0.5 x}, \quad x \ge 0$$
	\begin{enumerate}
		\item Calcule la probabilidad $P(X \le 1)$ de que el tiempo de espera sea menor a 1 ms.
		\item Calcule $P(2 \le X \le 5)$ mediante integración de la PDF.
		\item Halle la mediana $m$ tal que $P(X \le m) = 0.5$.
	\end{enumerate}
\end{problema}

\subsubsection*{Nivel Analítico}

\begin{problema}
	\label{prob:4.1.7}
	Sea $X$ una variable aleatoria continua con PDF $f(x)$. Demuestre formalmente que la esperanza matemática puede calcularse como:
	$$\E[X] = \int_{-\infty}^{\infty} x \cdot f(x)\,dx$$
	partiendo de la definición $\E[X] = \int x\,dF(x)$ y usando integración por partes bajo condiciones técnicas de regularidad.
\end{problema}

\begin{problema}
	\label{prob:4.1.8}
	Para una variable aleatoria continua $X$ con PDF $f(x)$ y una función $g: \mathbb{R} \to \mathbb{R}$, demuestre el \emph{Teorema LOTUS Continuo} (Law of the Unconscious Statistician):
	$$\E[g(X)] = \int_{-\infty}^{\infty} g(x) \cdot f(x)\,dx$$
	\begin{enumerate}
		\item Demuestre el caso particular $g(X) = X^{2}$.
		\item Derive la fórmula de la varianza $\Var(X) = \E[X^{2}] - (\E[X])^{2}$ usando LOTUS con $g(x) = (x - \mu)^{2}$.
	\end{enumerate}
\end{problema}

\subsubsection*{Nivel Desafiante}

\begin{problema}
	\label{prob:4.1.9}
	Considere una distribución mixta continua-discreta donde $X$ toma el valor $0$ con probabilidad $1/3$ y, condicionalmente a $X > 0$, sigue una distribución exponencial con parámetro $\lambda = 2$:
	$$f(x) = \begin{cases} 1/3 & \text{si } x = 0 \\ (2/3) \cdot 2 e^{-2x} & \text{si } x > 0 \\ 0 & \text{en otro caso} \end{cases}$$
	\begin{enumerate}
		\item Verifique que $f$ es una PDF válida (masa puntual en $0$ + densidad exponencial).
		\item Calcule la CDF $F(x)$ para todo $x \in \mathbb{R}$, identificando el salto en $x = 0$.
		\item Calcule $\E[X]$ y $\Var(X)$ usando la ley total de esperanza.
	\end{enumerate}
\end{problema}

\begin{problema}
	\label{prob:4.1.10}
	En teoría de señales, la convolución de dos densidades $f$ y $g$ produce la densidad de la suma $Z = X + Y$ de dos variables aleatorias independientes:
	$$h(z) = (f * g)(z) = \int_{-\infty}^{\infty} f(x) \cdot g(z - x)\,dx$$
	\begin{enumerate}
		\item Demuestre que si $X \sim N(0, 1)$ y $Y \sim N(0, 1)$ son independientes, entonces $Z = X + Y$ tiene PDF $N(0, 2)$, es decir, $\Var(Z) = \Var(X) + \Var(Y)$.
		\item Generalice: si $X_{1}, \dots, X_{n}$ son i.i.d. $N(0, \sigma^{2})$, entonces $\sum_{i=1}^{n} X_{i} \sim N(0, n\sigma^{2})$.
		\item Interprete este resultado en el contexto de filtrado de ruido gaussiano en telecomunicaciones.
	\end{enumerate}
\end{problema}

\subsection*{Sugerencias de los problemas --- Sección 04.01}

\begin{sugerencia}
	\label{sug:4.1.1}
	Para el Problema \ref{prob:4.1.1}, integre $\int_{0}^{\infty} c \cdot e^{-x}\,dx = c = 1$, así que $c = 1$. Para $P(1 \le X \le 3) = e^{-1} - e^{-3} \approx 0.318$. La CDF es $F(x) = 1 - e^{-x}$ para $x \ge 0$.
\end{sugerencia}

\begin{sugerencia}
	\label{sug:4.1.2}
	Para el Problema \ref{prob:4.1.2}, integre $\int_{0}^{2} c(4x - 2x^{2})\,dx = c \cdot [2x^{2} - 2x^{3}/3]_{0}^{2} = c(8 - 16/3) = 8c/3 = 1$, dando $c = 3/8$. Para $P(X \ge 1)$, integre desde $1$ hasta $2$.
\end{sugerencia}

\begin{sugerencia}
	\label{sug:4.1.3}
	Para el Problema \ref{prob:4.1.3}, la moda está en $x = 1$ (donde la PDF alcanza el máximo de $1$). La integral $\int_{0}^{1} x\,dx = 1/2$ y $\int_{1}^{2} (2-x)\,dx = 1/2$, sumando a $1$. La mediana $m$ satisface $\int_{0}^{m} x\,dx = 0.5$ si $m \le 1$, dando $m = 1$.
\end{sugerencia}

\begin{sugerencia}
	\label{sug:4.1.4}
	Para el Problema \ref{prob:4.1.4}, $f(x) = F'(x) = x e^{-x^{2}/2}$ para $x \ge 0$. Con la sustitución $u = x^{2}/2$, $du = x\,dx$: $\int_{0}^{\infty} x e^{-x^{2}/2}\,dx = \int_{0}^{\infty} e^{-u}\,du = 1$. La mediana satisface $1 - e^{-m^{2}/2} = 0.5$, dando $m = \sqrt{2 \ln 2} \approx 1.177$.
\end{sugerencia}

\begin{sugerencia}
	\label{sug:4.1.5}
	Para el Problema \ref{prob:4.1.5}, $\int_{-1}^{1} (3/8)(2 - x^{2})\,dx = (3/8)[2x - x^{3}/3]_{-1}^{1} = (3/8)(4 - 2/3) = (3/8)(10/3) = 10/8 \ne 1$. Verifique: $\int_{-1}^{1} (2 - x^{2})\,dx = [2x - x^{3}/3]_{-1}^{1} = (2 - 1/3) - (-2 + 1/3) = 4 - 2/3 = 10/3$. Por lo tanto $c = 3/10$ para que la integral sea $1$.
\end{sugerencia}

\begin{sugerencia}
	\label{sug:4.1.6}
	Para el Problema \ref{prob:4.1.6}, la CDF exponencial es $F(x) = 1 - e^{-0.5 x}$. Entonces $P(X \le 1) = 1 - e^{-0.5} \approx 0.393$. $P(2 \le X \le 5) = e^{-1} - e^{-2.5} \approx 0.3679 - 0.0821 = 0.286$. La mediana es $m = -\ln(0.5)/0.5 = 2 \ln 2 \approx 1.386$ ms.
\end{sugerencia}

\begin{sugerencia}
	\label{sug:4.1.7}
	Para el Problema \ref{prob:4.1.7}, use la representación $F(x) = \int_{-\infty}^{x} f(t)\,dt$ y la integración por partes $\int x\,dF(x) = xF(x) - \int F(x)\,dx$. Asumiendo $\lim_{x \to \pm\infty} xF(x) = 0$ (condición técnica de colas ligeras), sustituya la CDF y simplifique para obtener $\E[X] = \int x f(x)\,dx$.
\end{sugerencia}

\begin{sugerencia}
	\label{sug:4.1.8}
	Para el Problema \ref{prob:4.1.8}, use que para variables aleatorias continuas, $\E[g(X)] = \int g(x) f(x)\,dx$ donde $f$ es la densidad. El caso $g(X) = X^{2}$ se obtiene directamente. Para la varianza, $g(X) = (X - \mu)^{2}$ con $\mu = \E[X]$, dando $\E[(X-\mu)^{2}] = \E[X^{2}] - 2\mu\E[X] + \mu^{2} = \E[X^{2}] - \mu^{2}$.
\end{sugerencia}

\begin{sugerencia}
	\label{sug:4.1.9}
	Para el Problema \ref{prob:4.1.9}, la masa en $0$ es $1/3$ y la densidad exponencial reescalada $(2/3) \cdot 2 e^{-2x}$ integra a $\int_{0}^{\infty} (4/3) e^{-2x}\,dx = (4/3) \cdot (1/2) = 2/3$. Sumando: $1/3 + 2/3 = 1$. La CDF tiene un salto de $1/3$ en $x = 0$ y luego crece según la exponencial reescalada.
\end{sugerencia}

\begin{sugerencia}
	\label{sug:4.1.10}
	Para el Problema \ref{prob:4.1.10}, la convolución de dos gaussianas es gaussiana con varianzas sumadas. Demuestre usando la forma explícita de la PDF: $h(z) = \int (2\pi)^{-1} \exp(-x^{2}/2) \cdot \exp(-(z-x)^{2}/2)\,dx$. Complete el cuadrado en el exponente y use la integral gaussiana.
\end{sugerencia}

\subsection*{Soluciones de los problemas --- Sección 04.01}

\begin{solucion}
	\label{sol:4.1.1}
	Para el Problema \ref{prob:4.1.1}:
	\begin{enumerate}
		\item La condición de normalización requiere:
		\begin{align*}
			\int_{-\infty}^{\infty} f(x)\,dx = \int_{0}^{\infty} c \cdot e^{-x}\,dx = c \cdot [-e^{-x}]_{0}^{\infty} = c \cdot 1 = c = 1.
		\end{align*}
		Por lo tanto, $c = 1$.
		\item La probabilidad es:
		\begin{align*}
			P(1 \le X \le 3) = \int_{1}^{3} e^{-x}\,dx = [-e^{-x}]_{1}^{3} = e^{-1} - e^{-3} \approx 0.3679 - 0.0498 = 0.3181.
		\end{align*}
		\item La CDF se obtiene integrando desde $-\infty$ hasta $x$:
		\begin{align*}
			F(x) &= \int_{-\infty}^{x} f(t)\,dt = \begin{cases} 0 & \text{si } x < 0 \\ \int_{0}^{x} e^{-t}\,dt = 1 - e^{-x} & \text{si } x \ge 0 \end{cases}
		\end{align*}
		Esta es la CDF de la distribución exponencial con parámetro $\lambda = 1$.
	\end{enumerate}
\end{solucion}

\begin{solucion}
	\label{sol:4.1.2}
	Para el Problema \ref{prob:4.1.2}:
	\begin{enumerate}
		\item La condición de normalización es:
		\begin{align*}
			1 = \int_{0}^{2} c(4x - 2x^{2})\,dx = c \left[ 2x^{2} - \dfrac{2x^{3}}{3} \right]_{0}^{2} = c \left( 8 - \dfrac{16}{3} \right) = \dfrac{8c}{3}.
		\end{align*}
		Por lo tanto, $c = 3/8$.
		\item Para $P(X \ge 1)$:
		\begin{align*}
			P(X \ge 1) = \int_{1}^{2} \dfrac{3}{8}(4x - 2x^{2})\,dx = \dfrac{3}{8} \left[ 2x^{2} - \dfrac{2x^{3}}{3} \right]_{1}^{2} = \dfrac{3}{8} \cdot \left[ \left( 8 - \dfrac{16}{3} \right) - \left( 2 - \dfrac{2}{3} \right) \right] = \dfrac{3}{8} \cdot \dfrac{4}{3} = \dfrac{1}{2}.
		\end{align*}
		Para $P(0.5 \le X \le 1.5)$:
		\begin{align*}
			P(0.5 \le X \le 1.5) = \dfrac{3}{8} \left[ 2x^{2} - \dfrac{2x^{3}}{3} \right]_{0.5}^{1.5} = \dfrac{3}{8} \left[ \left( \dfrac{9}{2} - \dfrac{9}{4} \right) - \left( \dfrac{1}{2} - \dfrac{1}{12} \right) \right] = \dfrac{3}{8} \cdot \dfrac{5}{3} = \dfrac{5}{8}.
		\end{align*}
	\end{enumerate}
\end{solucion}

\begin{solucion}
	\label{sol:4.1.3}
	Para el Problema \ref{prob:4.1.3}:
	\begin{enumerate}
		\item Verificamos la normalización:
		\begin{align*}
			\int_{0}^{1} x\,dx + \int_{1}^{2} (2 - x)\,dx = \dfrac{1}{2} + \left[ 2x - \dfrac{x^{2}}{2} \right]_{1}^{2} = \dfrac{1}{2} + \left( 4 - 2 \right) - \left( 2 - \dfrac{1}{2} \right) = \dfrac{1}{2} + \dfrac{1}{2} = 1. \quad \checkmark
		\end{align*}
		\item La moda está en $x = 1$ (donde la PDF alcanza el máximo $f(1) = 1$). Para la mediana $m$: si $m \le 1$, $\int_{0}^{m} x\,dx = m^{2}/2 = 0.5$, dando $m = 1$. Como $m = 1$ también marca la frontera entre las dos ramas, es exactamente la mediana. Por la simetría de la triangular alrededor de $x = 1$, concluimos que \emph{moda $=$ mediana $= 1$}.
	\end{enumerate}
\end{solucion}

\begin{solucion}
	\label{sol:4.1.4}
	Para el Problema \ref{prob:4.1.4}:
	\begin{enumerate}
		\item Derivando la CDF:
		\begin{align*}
			f(x) = F'(x) = \dfrac{d}{dx}\left( 1 - e^{-x^{2}/2} \right) = x e^{-x^{2}/2}, \quad x \ge 0.
		\end{align*}
		\item Verificamos la normalización con la sustitución $u = x^{2}/2$, $du = x\,dx$:
		\begin{align*}
			\int_{0}^{\infty} x e^{-x^{2}/2}\,dx = \int_{0}^{\infty} e^{-u}\,du = [-e^{-u}]_{0}^{\infty} = 1 - 0 = 1. \quad \blacksquare
		\end{align*}
		\item La mediana satisface $1 - e^{-m^{2}/2} = 0.5$, es decir, $e^{-m^{2}/2} = 0.5$. Tomando logaritmo: $-m^{2}/2 = -\ln 2$, así $m = \sqrt{2 \ln 2} \approx 1.1774$.
	\end{enumerate}
\end{solucion}

\begin{solucion}
	\label{sol:4.1.5}
	Para el Problema \ref{prob:4.1.5}:
	\begin{enumerate}
		\item Verificamos la normalización:
		\begin{align*}
			\int_{-1}^{1} c(2 - x^{2})\,dx = c \left[ 2x - \dfrac{x^{3}}{3} \right]_{-1}^{1} = c \left[ \left( 2 - \dfrac{1}{3} \right) - \left( -2 + \dfrac{1}{3} \right) \right] = c \cdot \dfrac{10}{3}.
		\end{align*}
		Para que la integral sea $1$: $c = 3/10$. La PDF correcta es $f(x) = \dfrac{3}{10}(2 - x^{2})$ para $x \in [-1, 1]$.
		\item Para $P(-0.5 \le X \le 0.5)$:
		\begin{align*}
			P(-0.5 \le X \le 0.5) = \int_{-0.5}^{0.5} \dfrac{3}{10}(2 - x^{2})\,dx = \dfrac{3}{10}\left[ 2x - \dfrac{x^{3}}{3} \right]_{-0.5}^{0.5} = \dfrac{3}{10} \cdot 2 = 0.6.
		\end{align*}
		Para $P(X \ge 0)$ (por simetría de la PDF): $P(X \ge 0) = 0.5$.
		\item La CDF es:
		\begin{align*}
			F(x) = \int_{-1}^{x} \dfrac{3}{10}(2 - t^{2})\,dt = \dfrac{3}{10}\left( 2x + 2 - \dfrac{x^{3}}{3} - \dfrac{1}{3} \right) = \dfrac{3}{10}\left( \dfrac{6x - x^{3} + 5}{3} \right) = \dfrac{6x - x^{3} + 5}{10}.
		\end{align*}
		Verificación: $F(-1) = (-6 + 1 + 5)/10 = 0$ ✓, $F(1) = (6 - 1 + 5)/10 = 1$ ✓.
	\end{enumerate}
\end{solucion}

\begin{solucion}
	\label{sol:4.1.6}
	Para el Problema \ref{prob:4.1.6}, la CDF exponencial con parámetro $\lambda = 0.5$ es $F(x) = 1 - e^{-0.5 x}$ para $x \ge 0$.
	\begin{enumerate}
		\item $P(X \le 1) = 1 - e^{-0.5 \cdot 1} = 1 - e^{-0.5} \approx 1 - 0.6065 = 0.3935$.
		\item $P(2 \le X \le 5) = F(5) - F(2) = (1 - e^{-2.5}) - (1 - e^{-1}) = e^{-1} - e^{-2.5} \approx 0.3679 - 0.0821 = 0.2858$.
		\item La mediana satisface $1 - e^{-0.5 m} = 0.5$, es decir, $e^{-0.5 m} = 0.5$. Tomando logaritmo: $m = -2 \ln(0.5) = 2 \ln 2 \approx 1.386$ ms.
	\end{enumerate}
\end{solucion}

\begin{solucion}
	\label{sol:4.1.7}
	Para el Problema \ref{prob:4.1.7}, partimos de la definición $\E[X] = \int_{-\infty}^{\infty} x\,dF(x)$. Usando integración por partes con $u = x$ y $dv = dF(x)$:
	\begin{align*}
		\E[X] &= \int_{-\infty}^{\infty} x\,dF(x) = \left[ xF(x) \right]_{-\infty}^{\infty} - \int_{-\infty}^{\infty} F(x)\,dx.
	\end{align*}
	Asumiendo la condición técnica de colas ligeras $\lim_{x \to \pm\infty} xF(x) = 0$ (lo cual se cumple si $\int |x| f(x)\,dx < \infty$):
	\begin{align*}
		\E[X] &= - \int_{-\infty}^{\infty} F(x)\,dx.
	\end{align*}
	Sustituyendo $F(x) = \int_{-\infty}^{x} f(t)\,dt$:
	\begin{align*}
		\E[X] &= - \int_{-\infty}^{\infty} \int_{-\infty}^{x} f(t)\,dt\,dx = - \int_{-\infty}^{\infty} f(t) \left( \int_{t}^{\infty} dx \right) dt = - \int_{-\infty}^{\infty} f(t) \cdot (\infty - t)\,dt.
		\end{align*}
		Este enfoque no es el más limpio. Procedamos de manera más directa. Recordando que $dF(x) = f(x)\,dx$:
		\begin{align*}
			\E[X] = \int_{-\infty}^{\infty} x \cdot dF(x) = \int_{-\infty}^{\infty} x \cdot f(x)\,dx. \quad \blacksquare
		\end{align*}
		La sustitución directa $dF = f\,dx$ es la forma operacional del teorema, válida bajo la condición de integrabilidad absoluta $\int |x| f(x)\,dx < \infty$.
\end{solucion}

\begin{solucion}
	\label{sol:4.1.8}
	Para el Problema \ref{prob:4.1.8}:
	\begin{enumerate}
		\item Para $g(X) = X^{2}$, usamos que la PDF de $X^{2}$ se obtiene mediante la transformación. Sin embargo, la forma más directa es usar la definición operacional de LOTUS:
		\begin{align*}
			\E[X^{2}] = \int_{-\infty}^{\infty} x^{2} f(x)\,dx,
		\end{align*}
		que es simplemente la definición de esperanza con $g(x) = x^{2}$.
		\item Para la varianza con $g(X) = (X - \mu)^{2}$ donde $\mu = \E[X]$:
		\begin{align*}
			\Var(X) &= \E[(X - \mu)^{2}] = \int_{-\infty}^{\infty} (x - \mu)^{2} f(x)\,dx \\
			&= \int_{-\infty}^{\infty} (x^{2} - 2\mu x + \mu^{2}) f(x)\,dx \\
			&= \E[X^{2}] - 2\mu \E[X] + \mu^{2} \\
			&= \E[X^{2}] - 2\mu^{2} + \mu^{2} = \E[X^{2}] - (\E[X])^{2}. \quad \blacksquare
		\end{align*}
		donde usamos linealidad de la integral y la constancia de $\mu$ bajo la integral.
	\end{enumerate}
\end{solucion}

\begin{solucion}
	\label{sol:4.1.9}
	Para el Problema \ref{prob:4.1.9}:
	\begin{enumerate}
		\item La masa puntual en $x = 0$ es $1/3$, y la densidad exponencial reescalada $(2/3) \cdot 2 e^{-2x} = (4/3) e^{-2x}$ para $x > 0$ integra a:
		\begin{align*}
			\int_{0}^{\infty} \dfrac{4}{3} e^{-2x}\,dx = \dfrac{4}{3} \cdot \dfrac{1}{2} = \dfrac{2}{3}.
		\end{align*}
		Sumando: $1/3 + 2/3 = 1$. La PDF es válida.
		\item La CDF tiene un salto en $x = 0$ y luego crece continuamente:
		\begin{align*}
			F(x) = \begin{cases} 0 & \text{si } x < 0 \\ 1/3 & \text{si } x = 0 \\ 1/3 + \int_{0}^{x} \dfrac{4}{3} e^{-2t}\,dt = 1/3 + \dfrac{2}{3}(1 - e^{-2x}) = 1 - \dfrac{2}{3} e^{-2x} & \text{si } x > 0 \end{cases}
		\end{align*}
		El salto en $x = 0$ es $F(0) - F(0^{-}) = 1/3 - 0 = 1/3$, correspondiente a la masa puntual.
		\item Usando la ley total de esperanza $\E[X] = \E[\E[X \mid Z]]$ con $Z$ indicador de $X = 0$:
		\begin{align*}
			\E[X] = \Prob(X = 0) \cdot 0 + \Prob(X > 0) \cdot \E[X \mid X > 0].
		\end{align*}
		Condicionalmente a $X > 0$, $X$ sigue $\text{Exp}(2)$ con media $1/2$. Por lo tanto:
		\begin{align*}
			\E[X] = \dfrac{1}{3} \cdot 0 + \dfrac{2}{3} \cdot \dfrac{1}{2} = \dfrac{1}{3}.
		\end{align*}
		Para la varianza, $\Var(X) = \E[X^{2}] - (\E[X])^{2}$. Con $\E[X^{2} \mid X > 0] = 2/4 = 1/2$ (segundo momento de Exp(2)):
		\begin{align*}
			\E[X^{2}] = \dfrac{1}{3} \cdot 0 + \dfrac{2}{3} \cdot \dfrac{1}{2} = \dfrac{1}{3}, \quad \Var(X) = \dfrac{1}{3} - \dfrac{1}{9} = \dfrac{2}{9}. \quad \blacksquare
		\end{align*}
	\end{enumerate}
\end{solucion}

\begin{solucion}
	\label{sol:4.1.10}
	Para el Problema \ref{prob:4.1.10}:
	\begin{enumerate}
		\item La convolución de dos gaussianas estándar independientes $X, Y \sim N(0, 1)$ es:
		\begin{align*}
			h(z) &= (f * g)(z) = \int_{-\infty}^{\infty} \dfrac{1}{\sqrt{2\pi}} e^{-x^{2}/2} \cdot \dfrac{1}{\sqrt{2\pi}} e^{-(z-x)^{2}/2}\,dx \\
			&= \dfrac{1}{2\pi} \int_{-\infty}^{\infty} \exp\left( -\dfrac{x^{2} + (z-x)^{2}}{2} \right) dx.
		\end{align*}
		Expandiendo y completando el cuadrado en el exponente:
		\begin{align*}
			x^{2} + (z-x)^{2} &= 2x^{2} - 2zx + z^{2} = 2\left( x - \dfrac{z}{2} \right)^{2} + \dfrac{z^{2}}{2}.
		\end{align*}
		Por lo tanto:
		\begin{align*}
			h(z) &= \dfrac{1}{2\pi} e^{-z^{2}/4} \int_{-\infty}^{\infty} \exp\left( -\left( x - \dfrac{z}{2} \right)^{2} \right) dx.
		\end{align*}
		Con la sustitución $u = \sqrt{2}(x - z/2)$, $du = \sqrt{2}\,dx$:
		\begin{align*}
			\int_{-\infty}^{\infty} \exp\left( -\left( x - \dfrac{z}{2} \right)^{2} \right) dx = \sqrt{\pi}.
		\end{align*}
		Sustituyendo:
		\begin{align*}
			h(z) = \dfrac{1}{2\pi} \cdot \sqrt{\pi} \cdot e^{-z^{2}/4} = \dfrac{1}{\sqrt{2\pi} \cdot \sqrt{2}} e^{-z^{2}/(2 \cdot 2)} = \dfrac{1}{\sqrt{2\pi \cdot 2}} e^{-z^{2}/4}.
		\end{align*}
		Esta es la PDF de $N(0, 2)$. Por lo tanto, $Z = X + Y \sim N(0, 2)$. \quad \qedhere
		\item Por inducción: si $S_{n-1} = \sum_{i=1}^{n-1} X_{i} \sim N(0, (n-1)\sigma^{2})$ y $X_{n} \sim N(0, \sigma^{2})$ son independientes, entonces $S_{n} = S_{n-1} + X_{n}$ tiene convolución de dos gaussianas con varianzas $(n-1)\sigma^{2}$ y $\sigma^{2}$, dando $N(0, n\sigma^{2})$. La base inductiva $n = 1$ es trivial ($S_{1} = X_{1}$).
		\item En telecomunicaciones, si una señal $s(t)$ se transmite a través de un canal con ruido gaussiano aditivo $W(t) \sim N(0, \sigma^{2})$ y se promedian $n$ muestras independientes, la suma de los ruidos $W_{1} + \cdots + W_{n} \sim N(0, n\sigma^{2})$. La varianza de la suma crece linealmente con $n$, pero al dividir por $n$ para obtener el promedio, $\bar{W} \sim N(0, \sigma^{2}/n)$, cuya varianza decrece como $1/n$. Este es el principio fundamental del filtrado por promedio para reducción de ruido gaussiano blanco: promediar $n$ muestras reduce la potencia del ruido por un factor $n$.
	\end{enumerate}
\end{solucion}
